% Part: lambda-calculus
% Chapter: introduction
% Section: lambda-computable

\documentclass[../../../include/open-logic-section]{subfiles}

\begin{document}

\olfileid[id]{lam}{int}{cmp}
\olsection{Fungsi yang \usetoken{S}{lambda definable} Dapat Dihitung}

\begin{thm}
\ollabel{thm:lambda-computable}
Jika suatu fungsi parsial~$f$ !!{lambda defined} oleh suatu suku lambda, maka
fungsi itu dapat dihitung.
\end{thm}

\begin{proof}
Andaikan suatu fungsi~$f$ !!{lambda defined} oleh suku lambda~$X$. Mari kita
uraikan prosedur informal untuk menghitung~$f$. Pada masukan $m_0$,
\dots,~$m_{n-1}$, tuliskan suku $X \num m_0 \ldots \num m_{n-1}$. Bangun
sebuah pohon dengan mula-mula menuliskan semua reduksi satu-langkah dari suku
asal; di bawahnya, tuliskan semua reduksi satu-langkah dari suku-suku tersebut
(yakni reduksi dua-langkah dari suku asal); lalu teruskan proses ini. Jika suatu
numeral pernah tercapai, kembalikan numeral itu sebagai jawaban; jika tidak,
fungsi tersebut tidak terdefinisi.

Dengan menggunakan tesis Church, kita menyimpulkan bahwa fungsi ini dapat
dihitung. Cara yang lebih baik untuk membuktikan teorema ini ialah memberikan
deskripsi rekursif bagi prosedur pencarian tersebut. Sebagai contoh, kita dapat
mendefinisikan suatu barisan fungsi dan relasi rekursif primitif,
``$\fn{IsASubterm}$,'' ``$\fn{Substitute}$,''
``$\fn{ReducesToInOneStep}$,'' ``$\fn{ReductionSequence}$,''
``$\fn{Numeral}$,'' dan seterusnya. Prosedur rekursif parsial untuk menghitung
$f(m_0, \dots, m_{n-1})$ kemudian mencari barisan reduksi satu-langkah yang
dimulai dengan $X \num{m_0} \dots \num{m_{n-1}}$ dan berakhir dengan sebuah
numeral, lalu mengembalikan bilangan yang bersesuaian dengan numeral tersebut.
Perinciannya panjang dan menjemukan, tetapi selebihnya rutin.
\end{proof}

\end{document}
